\chapter{User documentation}

\section{What is Sudoku}

In an $n \times n$ Sudoku, place a number from $1$ to $n$ in every empty cell
such that each region contains every number exactly once.
Sudoku variants add more rules.
For more detailed explanation of all the rules, read the Chapter~\ref{ch:chap-1}.

\section{Creating a Sudoku Puzzle}

\subsection*{New Sudoku}

Before you can start solving a Sudoku puzzle, you have to start by creating the Sudoku task.
First step is to click the button "New Sudoku" and specify the size of the Sudoku, which can be from 3 to 16.

After you create the Sudoku task, press the button "Start Solving".

\subsection*{Regions}

According to the rules of Sudoku, every number has to be in each region exactly once.
Each row and column is a region by default, but you might add up to 3 more types of regions.
To add more regions, click the button "Edit Regions".

\subsubsection*{Rectangular/Jigsaw Regions}

To add Jigsaw regions, press the button "Jigsaw".
Hold the left mouse button to choose cells in a new Jigsaw region.
Right-click a Jigsaw region to remove it.
In an $n \times n$ Sudoku puzzle with Jigsaw regions,
there has to be exactly $n$ Jigsaw regions and each region has to have exactly $n$ cells.
All cells in a Jigsaw region must be edge-adjacent and no cell can be in two different Jigsaw regions at once.

If the size of the Sudoku puzzle is not a prime number, you can press the button "Rectangular",
which will let you choose any size of regions and the software will automatically add all of them.
For example $16 \times 16$ Sudoku can have Rectangular regions $2 \times 8$, $4 \times 4$, or $8 \times 2$.

To remove all Rectangular or Jigsaw regions, press the button "None".

\subsubsection*{Diagonal Regions}

To add or remove diagonal regions, press the button "Diagonal".

\subsubsection*{Extra Regions}

To add Extra regions, click the button "Custom".
Hold the left mouse button to choose cells in a new Extra region.
Right-click an Extra region to remove it.
In an $n \times n$ Sudoku puzzle with Extra regions, each Extra region has to have exactly $n$ cells, but
unlike the Jigsaw regions, there can be any amount of Extra regions.
All cells in an Extra region must be edge-adjacent and no cell can be in two different Extra regions at once.

If the size of the Sudoku puzzle is $4 \times 4$, $9 \times 9$, or $16 \times 16$,
you can press the "Windoku" button, which adds the specific pattern of Extra regions for Windoku Sudoku.

\subsection*{Numbers}

To add numbers, press the button "Edit Numbers".
To add a single number in a cell, left click it and press a number on the keyboard.
An empty cell can contain any number in the solution.
If you want to specify what combination of number can be in a cell, right click it,
then you can add or remove individual numbers.
Press the \texttt{Delete} or the \texttt{Backspace} to make a cell empty again.
In a $10 \times 10$ or bigger Sudoku puzzle, you can press the "+10" button to add numbers bigger than 9.

\subsection*{Variants}

To add a Sudoku variant, press the button "Edit Variant".
You can choose one of five different variants to add.
To add labels to the edges, left-click the square on the edge, which will toggle it to the next label.
To toggle labels in the opposite direction, right-click it instead.

\subsection*{Empty Labels}

.

\section{Solving a Sudoku}

After you create a Sudoku task, press the button "Start Solving".
The software will automatically fill every empty cell with all possible numbers.

The goal of solving is to remove numbers from cells using different strategies,
until there is only one number in each cell.
For list and explanation of all the strategies used by the software, read the Chapter~\ref{ch:chap-2}.

The available strategies are ordered from the most simple to the most complicated.
If you want to turn off a strategy, press its button in the right menu.

To solve the Sudoku puzzle, press the button "Apply Step".
The software will go through all the strategies from the most simple to the most complicated and
find the first one that can prove that some of the numbers can not be a part of the solution.
Then it will show you the explanation of the strategy.
To continue, press the button "Apply Step" again, and the software will remove all numbers with yellow background.
For explanation of the colors that the software uses, read the first part of Chapter~\ref{ch:chap-2}.
(TODO: make it a section)

You can edit the numbers manually.
To set a cell to have a single number, left-click it.
To add or remove a single number from a cell, right-click it.
Press the \texttt{Delete} or the \texttt{Backspace} to reset a cell to what it was in the task.
In a $10 \times 10$ or bigger Sudoku puzzle, you can press the "+10" button to add or remove numbers bigger than 9.

If the task has no solution, or you have manually removed a number that has to be in the solution,
you can end up with some cells not being able to have any numbers.
If there is a cell with "!", it meas that the cell can not have any number.

\section{Examples of Sudoku Puzzles}

Visit Chapter~\ref{ch:chap-1} for examples of Sudoku tasks with unique solution.
